% !PRINTED 380
\weberpage{380}
\S\,104. \gesperrt{Charaktere und Geschlechter der quadratischen Formen.}

Unter einem Stammteiler $\delta$ einer Diskriminante $D$ wollen
wir folgendes verstehen (\S\,84):

1. \gesperrt{$\delta$ ist eine in $D$ aufgehende Stammdiskriminante.}

2. \gesperrt{Der Quotient $D:\delta = D_1$ ist selbst noch Diskriminante.}

Ein ungerader Stammteiler kann keine anderen Primfaktoren
enthalten als solche, die in $D$ aufgehen, und keinen mehr als
einmal; dagegen ist ein beliebiges Produkt aus verschiedenen in
$D$ aufgehenden ungeraden Primzahlen, mit einem solchen Vorzeichen
 versehen, daß $\delta \equiv 1 \pmod{4}$ wird, immer ein Stammteiler.

Außerdem kommt noch unter den Stammteilern vor:

\weq{1}{\begin{aligned}
&-4, \text{ wenn } D \equiv 0,\ -4 \pmod{16}, \\
&+8, \text{ \glqq } D \equiv 0,\ 8 \pmod{32}, \\
&-8, \text{ \glqq } D \equiv 0,\ -8 \pmod{32},
\end{aligned}}

Bezeichnen wir also mit $\delta_1$ die ungeraden Stammteiler, so
ergeben sich die sämtlichen Stammteiler $\delta$:
% [sic] Weber reuses the symbol delta_1 here for the odd stem-divisors, distinct from the delta_1 of the abstract-group discussion further below on this page.

\[
\begin{aligned}
1)\ & \delta = \delta_1, & D &\equiv\ \ \ 1 \pmod{4}, \\
& & &\equiv\ \ \ 4 \pmod{16}, \\
2)\ & \delta = \delta_1,\ -4\delta_1, & D &\equiv -4 \pmod{16}, \\
& & &\equiv\ 16 \pmod{32}, \\
3)\ & \delta = \delta_1,\ +8\delta_1, & D &\equiv\ \ \ 8 \pmod{32}, \\
4)\ & \delta = \delta_1,\ -8\delta_1, & D &\equiv -8 \pmod{32}, \\
5)\ & \delta = \delta_1,\ -4\delta_1, & D &\equiv\ \ \ 0 \pmod{32}. \\
& \ \ \ \ 8\delta_1,\ -8\delta_1, & &
\end{aligned}
\]

Die Anzahl der Stammteiler (1 als Stammteiler mitgerechnet)
ist hiernach stets eine Potenz von 2.

3. \gesperrt{Setzen wir sie gleich $2^\lambda$, so ist $\lambda$ in den Fällen 1), 2),
3), 4) gleich der Anzahl der in $D$ aufgehenden verschiedenen
Primzahlen, in dem Falle 5) um eins größer.}

Die Stammteiler von $D$ lassen sich durch Anwendung einer
symbolischen Multiplikation zu einer \gesperrt{Abelschen Gruppe} machen.

Wenn nämlich $\delta_1$, $\delta_2$ zwei Stammteiler von $D$ sind, so ist
das Produkt $\delta_1\delta_2$ auch eine Diskriminante, aber nicht immer
eine Stammdiskriminante. Wir bezeichnen mit $\delta$ den \gesperrt{Stamm}
von $\delta_1\delta_2$ und setzen symbolisch

\weq{2}{\delta = \delta_1\delta_2.}
